\section{第11章\quad 曲线积分}
\subsection{对弧长的曲线积分}

\begin{exercise}{392}{弧长积分与取向}
设分段光滑曲线弧 $L$ 的长度为 $a>0$，$k$ 为常数。从弧长积分的定义求 $\int_Lk\,\dd s$，并判断将曲线反向后结果是否改变；若已知 $\int_L(k+1)\,\dd s=3a$，求 $k$。
\end{exercise}

\begin{exercise}{393}{弧长积分的保序性}
设 $L$ 为长度非零的分段光滑曲线弧，$f$ 在 $L$ 上连续且非负。证明 $\int_Lf\,\dd s\geq0$；进一步证明，若 $f$ 在 $L$ 的某个内点处为正，则积分严格为正，并据此给出积分等于零的充要条件。
\end{exercise}

\begin{exercise}{394}{积分比较与严格不等式}
设 $f,g$ 在长度非零的分段光滑曲线弧 $L$ 上连续，且 $f\geq g$。证明
$\int_Lf\,\dd s\geq\int_Lg\,\dd s$。
判断仅有“某一点 $f>g$”能否推出严格不等式，并说明连续性在判断中的作用。
\end{exercise}

\begin{exercise}{395}{积分绝对值不等式的等号}
设实值连续函数 $f$ 定义在长度非零的空间光滑曲线弧 $L$ 上。证明
$|\int_Lf\,\dd s|\leq\int_L|f|\,\dd s$，并给出等号成立的充要条件；用单位圆上的 $f(x,y)=x$ 说明不等式可以严格成立。
\end{exercise}

\begin{exercise}{396}{圆弧上径向函数的积分}
设 $R>0$，$L_R$ 为下半圆 $y=-\sqrt{R^2-x^2}$。求 $\int_{L_R}\sqrt{x^2+y^2}\,\dd s$ 及径向距离的平均值，给出 $R=2$ 的结果，并解释两者随 $R$ 的缩放次数。
\end{exercise}

\begin{exercise}{397}{线段上的仿射函数平均值}
设 $L$ 连接 $(-1,0)$ 与 $(0,2)$。计算 $\int_L(2x-y+5)\,\dd s$；证明任意仿射函数在线段上的弧长平均值等于其在中点处的值，并用该结论检验结果。
\end{exercise}

\begin{exercise}{398}{水平圆周与参数积分}
设 $R>0,h\in\mathbb R$，$L$ 为 $x^2+y^2=R^2,z=h$ 的圆周。求 $\oint_L\arctan\frac{z^2}{x^2+y^2}\,\dd s$，给出 $R=h=2$ 的结果；固定 $R$ 时讨论积分随 $|h|$ 的单调性与极限。
\end{exercise}

\begin{exercise}{399}{大圆上的轮换对称性}
设 $\Gamma$ 为球面 $x^2+y^2+z^2=9$ 与平面 $x+y+z=0$ 的交圆。求 $\oint_\Gamma z^2\,\dd s$，再求 $\oint_\Gamma xy\,\dd s$；说明坐标轮换与平面方程分别提供了什么关系。
\end{exercise}

\begin{exercise}{400}{椭圆约束下的二次函数}
设椭圆 $L:x^2/4+y^2/3=1$ 的周长为 $a$。求所有使 $\lambda x^2+\mu y^2$ 在 $L$ 上为常数的 $(\lambda,\mu)$，据此求 $\oint_L(2xy+3x^2+4y^2)\,\dd s$。不要求求出周长 $a$。
\end{exercise}

\begin{exercise}{401}{折线对称性与积分比例}
设 $L_1:y=|x|$（$-1\leq x\leq1$），$L_2:y=x$（$0\leq x\leq1$）。分别计算
$I_1=\int_{L_1}(xy^2-|x|+2y)\,\dd s$ 与 $I_2=\int_{L_2}x\,\dd s$，求其比例；将 $2y$ 换成 $cy$ 后，求使新 $I_1=0$ 的 $c$。
\end{exercise}

\begin{exercise}{402}{线段上函数积分的排序}
设 $L$ 连接 $(1,0)$ 与 $(0,1)$，
$I=\int_Lx^2\,\dd s,J=\int_L(x-y)\,\dd s,K=\int_L(x+y)/2\,\dd s$。
求三者并严格排序；若 $I_p=\int_Lx^p\,\dd s$（$p>0$），求使 $J<I_p<K$ 成立的全部 $p$。
\end{exercise}

\begin{exercise}{403}{互换对称性与分母条件}
设 $0<R<(\sqrt3/2)\ln\pi$，$L:x^2+xy+y^2=R^2$ 的周长为 $a$。证明 $\sin(e^x)+\sin(e^y)>0$ 在 $L$ 上成立，求
$\oint_L\frac{2\sin(e^x)-3\sin(e^y)}{\sin(e^x)+\sin(e^y)}\,\dd s$。
说明去掉对 $R$ 的限制后，能否仅凭对称性就忽略积分存在性。
\end{exercise}

\begin{exercise}{404}{球面截圆上的交叉矩}
设 $L$ 为 $x^2+y^2+z^2=1$ 与 $x+y+z=0$ 的交圆。求 $\oint_L(xy+z)\,\dd s$，并求任意常数 $a,b,c$ 下 $\oint_L(ax y+byz+czx)\,\dd s$。解释线性项与交叉二次项所用对称性为何不同。
\end{exercise}

\begin{exercise}{405}{指数曲线与弧长因子}
设 $b>0$，$L_b:y=e^x$（$0\leq x\leq b$）。求 $\int_{L_b}\sqrt{1+y^2}\,\dd s$，给出 $b=1$ 的结果；判断其是否等于 $L_b$ 的长度并给出严格比较。
\end{exercise}

\begin{exercise}{406}{抛物线参数与弧长平均}
设 $b>0$，$L_b:y^2=2x$（$0\leq y\leq b$）。计算 $\int_{L_b}y\,\dd s$，给出 $b=\sqrt3$ 的结果；写出曲线长度的一元积分式，证明 $y$ 的弧长平均值位于 $(0,b)$。
\end{exercise}

\begin{exercise}{407}{双曲线弧上的因子约简}
设 $0<a<1$，$L_a$ 为 $xy=1$ 上 $(a,1/a)$ 至 $(1,1)$ 的弧段。计算 $\int_{L_a}\frac{x+y}{\sqrt{x^2+y^2}}\,\dd s$，给出 $a=1/2$ 的结果；用 $x$ 与用 $y$ 作参数分别写出积分，验证换向不改变结果。
\end{exercise}

\begin{exercise}{408}{平移圆弧与角度参数}
设 $a>0,0<b\leq\pi$，$L:x=a(1+\cos t),y=a\sin t$（$0\leq t\leq b$）。求 $\int_Lxy\,\dd s$，给出 $b=\pi/2$ 的结果；固定 $b$，说明积分关于 $a$ 的缩放次数。
\end{exercise}

\begin{exercise}{409}{极坐标闭曲线与几何中心}
设 $L$ 的极坐标方程为 $r=\cos\theta+\sin\theta$，$-\pi/4\leq\theta\leq3\pi/4$。求其直角坐标方程、圆心和半径，计算 $\oint_L(x+y)\,\dd s$；再求 $x+y$ 的弧长平均值并用圆心坐标解释。
\end{exercise}

\begin{exercise}{410}{空间螺旋弧与收敛极限}
设 $T>0$，$\Gamma_T:x=e^t\cos t,y=e^t\sin t,z=e^t$（$0\leq t\leq T$）。求
$\int_{\Gamma_T}(x^2+y^2+z^2)^{-1}\,\dd s$，给出 $T=2$ 的结果；讨论 $T\to+\infty$ 时曲线长度与该积分是否各有有限极限。
\end{exercise}

\begin{exercise}{411}{对数曲线的参数选择}
设 $b>1$，$L_b:y=\ln x$（$1\leq x\leq b$）。求 $\int_{L_b}\frac{e^y}{1+x^2}\,\dd s$，给出 $b=\sqrt3$ 的结果；证明该积分随 $b$ 严格增加，并求 $b\to+\infty$ 时的极限。
\end{exercise}

\begin{exercise}{412}{椭圆弧与根式抵消}
设 $L$ 为椭圆 $x^2+4y^2=2x+3$ 在 $x$ 轴上方的部分。选取参数使 $\sqrt{1+3y^2}\,\dd s$ 化简，计算其积分；若改取完整椭圆，结果如何？说明对称性依据。
\end{exercise}

\begin{exercise}{413}{双纽线的极坐标分段}
设 $a>0$，$L:(x^2+y^2)^2=a^2(x^2-y^2)$ 为完整双纽线。计算 $\oint_L(|x|-|y|)\,\dd s$，写明两叶曲线的角度范围，并从曲线方程证明被积函数处处非负。
\end{exercise}

\begin{exercise}{414}{空间线段的弧长积分}
设 $L$ 连接 $A(-1,0,2)$ 与 $B(3,-2,1)$。用同一个参数分别计算 $\int_Lxyz\,\dd s$ 与 $\int_L|xyz|\,\dd s$，明确第二个积分的分段点，并判断二者是否互为绝对值。
\end{exercise}

\begin{exercise}{415}{空间截圆上的二次矩}
设 $|c|<2\sqrt3$，$L_c$ 为 $x^2+y^2+z^2=6$ 与 $x+z=c$ 的交圆。求圆心、半径及 $\oint_{L_c}(xz+y^2)\,\dd s$，给出 $c=2$ 的结果；可取沿平面内两条正交直径的坐标简化二次项。
\end{exercise}

\begin{exercise}{416}{封闭边界与线密度质量}
设 $L$ 为直线 $y=2x$ 与抛物线 $y^2=x$ 围成区域的完整边界。求两交点，分段计算 $\oint_Ly\,\dd s$；把 $y$ 视为线密度时，求两段细丝的质量比例，并说明端点密度为零是否妨碍总质量为正。
\end{exercise}

\begin{exercise}{417}{折线上的弧长积分}
设 $L$ 为顶点 $O(0,0),A(1,0),B(0,1)$ 的三角形边界。分边计算
$\oint_L\frac{x}{1+y}\,\dd s$，并求密度为 $\rho(x,y)=1/(1+y)$ 的三角形细丝的质心横坐标。
\end{exercise}

\begin{exercise}{418}{弧长平均值与平方积分}
设 $L$ 为长度 $\ell>0$ 的光滑曲线弧，$f$ 为实值连续函数。证明
$|\int_Lf\,\dd s|^2\leq\ell\int_Lf^2\,\dd s$，给出等号成立的充要条件。
比较“$\int_Lf\,\dd s=0$”与“$\int_Lf^2\,\dd s=0$”分别能推出什么；用单位圆上的一个具体连续函数说明两条件不能互相替代。
\end{exercise}

\begin{exercise}{419}{圆形细丝的非均匀密度}
设 $L:x^2+y^2=1$，线密度为 $\rho_\lambda(x,y)=1+\lambda x$。
求使密度处处非负的全部 $\lambda$，在该范围内求质量、质心及关于 $y$ 轴的转动惯量。进一步求过质心且平行于 $y$ 轴的直线的惯量，讨论其随 $|\lambda|$ 的变化，并说明曲线反向是否影响各物理量。
\end{exercise}

\subsection{对坐标的曲线积分}

\begin{exercise}{420}{坐标微分与弧长微分}
设 $L$ 为从 $(c,a)$ 到 $(c,b)$ 的有向线段，其中 $a<b$，$P,Q$ 连续。求
$\int_LP(x,y)\,\dd x$ 和 $\int_LQ(x,y)\,\dd y$ 的一元积分表达式；再取 $P\equiv1$，比较 $\int_LP\,\dd x$ 与 $\int_LP\,\dd s$。
\end{exercise}

\begin{exercise}{421}{常向量场与端点}
设 $L$ 为由 $A(x_1,y_1)$ 到 $B(x_2,y_2)$ 的分段光滑曲线，$k,m$ 为常数。求
$\int_Lk\,\dd x$、$\int_Lk\,\dd y$ 及 $\int_L(k\,\dd x+m\,\dd y)$；据此说明常向量场所做的功与路径无关，并求反向积分。
\end{exercise}

\begin{exercise}{422}{三角根式与绝对值}
设 $0<b\leq\pi$，$L_b:y=\sin x$ 从 $(0,0)$ 到 $(b,\sin b)$。求 $I(b)=\int_{L_b}\sqrt{1-y^2}\,\dd x$ 的分段表达式，给出 $I(\pi)$，并指出直接以 $\cos x$ 代替根式会在哪个区间产生错误。
\end{exercise}

\begin{exercise}{423}{抛物线上的奇函数消去}
设 $a>0$，$L_a:y=x^2$ 从 $(-a,a^2)$ 到 $(a,a^2)$。计算 $\int_{L_a}(x+e^y)\,\dd y$，给出 $a=1$ 的结果；分别说明 $x$ 项与 $e^y$ 项是否为零，不能仅根据端点纵坐标相同判断整个积分。
\end{exercise}

\begin{exercise}{424}{线段参数与方向符号}
设 $a,b>0$，$L$ 为从 $(a,0)$ 到 $(0,b)$ 的线段。求 $\int_Lxy\,\dd y$ 与 $\int_Lxy\,\dd x$，给出 $a=3,b=1$ 时前者的值；解释两积分符号不同的原因。
\end{exercise}

\begin{exercise}{425}{双曲线上的对数积分}
设 $b>1$，$L$ 为 $xy=1$ 上从 $(1,1)$ 到 $(b,1/b)$ 的弧。求 $\int_Lx\ln y\,\dd y$，给出 $b=2$ 的结果；用起终点及 $\ln y$ 的符号独立判断结果的正负。
\end{exercise}

\begin{exercise}{426}{曲线约束与恰当微分}
设 $L:y=\cos x$ 从 $(0,1)$ 到 $(\pi,-1)$。计算 $\int_Le^y\sin x\,\dd x+(1-y)\,\dd y$；判断该微分形式在整个平面内是否与路径无关，解释本题沿给定曲线仍可快速求值的原因。
\end{exercise}

\begin{exercise}{427}{回折抛物线的参数选择}
设 $L:y^2=x$ 从 $(1,-1)$ 到 $(1,1)$。取 $y$ 为参数计算 $\int_L\frac{xy\,\dd x-x^2\,\dd y}{1+x}$；说明直接以 $x$ 为全程单调参数会遗漏什么，并写出以 $x$ 分段时的两段方向。
\end{exercise}

\begin{exercise}{428}{摆线一拱与端点项}
设 $a>0$，$L:x=a(t-\sin t),y=a(1-\cos t)$，$t:0\to2\pi$。计算 $\int_L(2a-y)\,\dd x+\dd y$，并求 $\int_Ly\,\dd x$；用两者之和与水平位移检验结果。
\end{exercise}

\begin{exercise}{429}{螺旋线上的多项式三角积分}
设 $T>0$，$L_T:x=2\cos t,y=2\sin t,z=t$，$t:0\to T$。求 $\int_{L_T}y\,\dd x+z\,\dd y-x^2\,\dd z$，给出 $T=\pi$ 的结果；另求将路径改为同端点直线段后的积分，判断是否与路径无关。
\end{exercise}

\begin{exercise}{430}{椭圆弧上的端点法}
设 $L$ 为椭圆 $x^2+4y^2=4$ 在第一象限内按逆时针方向的弧。找出 $y\,\dd x+x\,\dd y$ 的原函数，计算积分；若终点改为该弧上任意点 $(u,v)$，写出积分值并求其最大值。
\end{exercise}

\begin{exercise}{431}{角形式与绕行次数}
设 $a>0$，$L$ 为顺时针圆周 $x^2+y^2=a^2$。参数化计算 $\oint_L\frac{y\,\dd x-x\,\dd y}{x^2+y^2}$；若沿同一圆周顺时针走 $m$ 圈（$m$ 为正整数），结果如何？说明为何不能在整个圆盘上直接使用无奇点的格林公式。
\end{exercise}

\begin{exercise}{432}{圆周约束与分母退化}
设 $r>0$，$L_r:(x-1)^2+(y-2)^2=r^2$ 取逆时针方向。讨论 $\oint_{L_r}\frac{\dd x+\dd y}{x^2+y^2-2x-4y}$ 何时有定义，有定义时计算，给出 $r=3$ 的结果；特别说明 $r=\sqrt5$ 的情形。
\end{exercise}

\begin{exercise}{433}{空间径向场与路径替换}
设 $L$ 从 $A(1,2,3)$ 到 $B(-1,4,0)$。证明 $\int_Lx\,\dd x+y\,\dd y+z\,\dd z$ 与路径无关并计算；若终点可在球面 $x^2+y^2+z^2=17$ 上任意变化，积分是否改变？
\end{exercise}

\begin{exercise}{434}{斜截圆的观察方向}
设 $\Gamma:x^2+y^2=1,x+z=1$，从 $x$ 轴正方向观察为顺时针。写出符合方向的参数方程，计算
$\oint_\Gamma yz\,\dd x-\frac{x}{x^2+y^2}\,\dd y-\frac{y}{x+z}\,\dd z$。
分别核查两个分母在曲线上是否为零，并给出反向后的值。
\end{exercise}

\begin{exercise}{435}{折线拐点与分段积分}
设 $L:y=1-|x|$，从 $(-1,0)$ 到 $(1,0)$。计算 $\int_Ly^2\,\dd x+x^2\,\dd y$；把顶点 $(0,1)$ 改为 $(0,h)$（$h>0$）后求积分关于 $h$ 的表达式，说明两侧 $\dd y$ 的符号。
\end{exercise}

\begin{exercise}{436}{三角形边界上的指数场}
设 $h>0$，$L_h$ 为依次经过 $O(0,0),A(1,0),B(1,h),O$ 的闭折线。分边求 $\oint_{L_h}e^y\,\dd x+e^x\,\dd y$，给出 $h=2$ 的结果；说明不能只根据闭合性将积分判为零。
\end{exercise}

\begin{exercise}{437}{复合圆弧与全微分}
设 $L$ 从 $O(0,0)$ 沿第一象限的圆弧 $x^2+y^2=2x$ 到 $A(2,0)$，再沿 $x^2+y^2=4$ 的第一象限圆弧到 $B(0,2)$。求 $\int_Ly\,\dd x+(x-1)\,\dd y$，分别用分段法和原函数法计算；说明若中间点及路径改变但端点不变，结果是否保持。
\end{exercise}

\begin{exercise}{438}{角点的单侧极限与曲线积分}
设 $L$ 从原点沿 $y=\sqrt x$ 到 $(1,1)$，再沿 $y=x$ 回原点。分段计算 $\oint_L\arctan(x/y)\,\dd x-\dd y$，各段端点按单侧极限处理；判断被积函数能否在原点作连续延拓，说明仅改原点处的有限取值是否影响积分。
\end{exercise}

\begin{exercise}{439}{空间折线上的坐标微分}
设 $L$ 依次经过 $A(0,0,0),B(0,0,2),C(1,0,2),D(1,3,2)$。求 $\int_Lz^2\,\dd x+x^2\,\dd y+y^2\,\dd z$；再按 $A\to(1,0,0)\to(1,3,0)\to D$ 的顺序移动，求新积分并比较，指出哪一项体现了路径影响。
\end{exercise}

\begin{exercise}{440}{空间交圆的平面化}
设 $L:x^2+y^2+z^2=6,y+z=2$，从 $z$ 轴正方向观察为逆时针。利用 $y+z=2$ 将
$\oint_Lz\,\dd x+x\,\dd y-\frac{y+z}{2}\,\dd z$ 化为平面投影上的积分，求值并说明投影方向；再用一组正交参数检验符号。
\end{exercise}

\begin{exercise}{441}{折点路径与全微分}
设 $L$ 为圆锥面 $z=\sqrt{x^2+y^2}$ 与平面 $y=x$ 的交线上从 $A(1,1,\sqrt2)$ 经原点到 $B(-1,-1,\sqrt2)$ 的曲线弧。
说明 $L$ 在原点处的光滑性，并分别用分段参数法和全微分法计算
$\int_Lyz\,\dd x+zx\,\dd y+xy\,\dd z$。
\end{exercise}

\begin{exercise}{442}{有向单位切向量}
设 $L$ 是圆周 $x^2+y^2=2y$ 上从原点到 $(0,2)$ 的右半圆弧，$P,Q$ 连续。
求与 $L$ 同向的单位切向量，将 $\int_LP\,\dd x+Q\,\dd y$ 化为对弧长的积分；再取 $P=x,Q=y$ 计算，并检验反向后的符号。
\end{exercise}

\begin{exercise}{443}{两类曲线积分与方向}
设 $L$ 为 $x=t,y=t^2,z=t^3$ 上 $t$ 从 $1$ 变到 $0$ 的曲线弧，弧长为 $a$。
不用求出 $a$ 的具体数值，计算
\[
 \int_L\frac{\dd x+2x\,\dd y+3y\,\dd z}{\sqrt{1+4x^2+9y^2}}.
\]
再将三个坐标微分的系数解释为沿正参数方向的单位切向量分量，说明取向如何决定结果。
\end{exercise}

\begin{exercise}{444}{抛物线弧上的积分转换}
设 $f$ 连续，$L$ 为 $y=x^2$ 上从 $(1,1)$ 到 $(0,0)$ 的曲线弧。
\begin{enumerate}[label=(\arabic*)]
\item 把 $I=\int_L f(x,y)\,\dd x+2xf(x,y)\,\dd y$ 转化为对弧长的积分，写明符号的来源。
\item 当 $f(x,y)=1/(x^2+3y+1)$ 时计算 $I$。
\item 当 $f$ 在 $L$ 上非负时，判断 $I$ 的符号；给出等号成立的条件。
\end{enumerate}
\end{exercise}

\begin{exercise}{445}{非单调参数与重复经过的弧段}
设 $\gamma(u)=(u,u^2)$（$0\leq u\leq1$）。曲线 $C$ 由 $u=4t(1-t)$、$t:0\to1$ 描述，因此先沿抛物线从原点走到 $(1,1)$，再原路返回。
分别计算 $\int_Cx\,\dd s$ 与 $\int_C(-y\,\dd x+x\,\dd y)$，比较它们与只走去程的相应积分。解释第一类积分中的参数速度为何取绝对值，以及为何不能把 $u(0)=u(1)$ 直接代入第一类积分的积分限判为零。
\end{exercise}

\begin{exercise}{446}{路径族与保守场参数}
设 $\alpha\in\mathbb R,p\geq1$，$L_p:y=x^p$ 从 $(0,0)$ 到 $(1,1)$。求
$W_\alpha(p)=\int_{L_p}\alpha y\,\dd x+x\,\dd y$。
求使 $W_\alpha(p)$ 对所有 $p\geq1$ 均相同的 $\alpha$，证明相应向量场在整个平面内保守；对其余 $\alpha$ 讨论 $W_\alpha$ 关于 $p$ 的单调性与 $p\to+\infty$ 时的极限，并解释路径趋近折线时的结果。
\end{exercise}

\subsection{格林公式及其应用}
\subsubsection{封闭曲线、补线与奇点}

\begin{exercise}{447}{闭曲线积分与全微分}
设 $L$ 为任意分段光滑有向闭曲线。求 $\oint_L(y-\cos x)\,\dd x+(x+\sin y)\,\dd y$，并找出原函数；说明即使闭曲线自交，原函数法是否仍然适用。
\end{exercise}

\begin{exercise}{448}{环流数据与面积确定}
设 $D$ 为有界区域，边界分段光滑且取正向，已知 $\oint_{\partial D}(1+3y)\,\dd x+(5x-1)\,\dd y=4$。求 $D$ 的面积；把 $3,5$ 分别换为参数 $a,b$ 后，讨论何时可以由同一积分数据唯一确定面积，何时数据本身不可能成立。
\end{exercise}

\begin{exercise}{449}{负向三角形边界}
设 $a,b>0$，$L$ 为顶点 $(0,0),(a,0),(a,b)$ 的三角形边界，取顺时针方向。求 $\oint_L(2x+y-4)\,\dd x+(5x-3y+6)\,\dd y$，给出 $a=3,b=2$ 的结果；指出其中哪些项对任意闭曲线的积分均为零。
\end{exercise}

\begin{exercise}{450}{双曲线区域的负向环流}
设 $D$ 由 $xy=1,y=x,x=2$ 围成，$L$ 为其负向边界。计算 $\oint_L(2x+y^2)\,\dd x+(3y-x^2)\,\dd y$，并由格林公式中的被积函数及取向先判断符号；另写出正向边界的结果。
\end{exercise}

\begin{exercise}{451}{圆周环流与四次缩放}
设 $R>0$，$L$ 为正向圆周 $x^2+y^2=R^2$。求 $\oint_Lx^2y\,\dd x-xy^2\,\dd y$，并求积分与圆盘面积之比；分别从变量缩放与格林公式解释结果关于 $R$ 的次数。
\end{exercise}

\begin{exercise}{452}{扇环边界与反正切场}
设 $0<r<R$，$D=\{(\rho\cos\theta,\rho\sin\theta):r\leq\rho\leq R,\pi/4\leq\theta\leq\pi/3\}$。计算其正向边界上的
$\oint_{\partial D}\frac1x\arctan\frac yx\,\dd x+\frac2y\arctan\frac xy\,\dd y$，给出 $r=1,R=2$ 的结果；检查坐标轴上的奇点为何不妨碍使用格林公式。
\end{exercise}

\begin{exercise}{453}{抛物线弓形的面积与环流}
设 $L$ 为 $y^2=2x,y=x-4$ 围成区域的顺时针边界。求交点、区域面积和形心纵坐标，据此计算 $\oint_L(e^x+2y)\,\dd x-(2e^y-xy)\,\dd y$，说明指数项对闭合积分的贡献。
\end{exercise}

\begin{exercise}{454}{菱形域与奇函数消去}
设 $a>0$，$L_a$ 为菱形 $|x|+|y|\leq a$ 的逆时针边界。计算 $\oint_{L_a}(x^2y-xy^3)\,\dd x+(y^2-2xy)\,\dd y$，给出 $a=1$ 的结果；写出格林公式转换后各项被保留或消去的依据。
\end{exercise}

\begin{exercise}{455}{补线与对称抵消}
设 $h>0$，$L_h$ 依次连接 $(-1,0),(0,h),(1,0)$。用补线法求 $\int_{L_h}xy\,\dd x+x^2\,\dd y$，给出 $h=1$ 的结果；再用两段参数验证，并解释其与 $h$ 无关的原因。
\end{exercise}

\begin{exercise}{456}{半圆弧补线与端点项}
设 $L:y=\sqrt{4x-x^2}$ 从 $O(0,0)$ 到 $A(4,0)$。计算 $\int_L(x^2-3y)\,\dd x+(2y^3+x)\,\dd y$，分别标明原函数项与面积项的贡献；若改走下半圆且端点顺序不变，两结果相差多少？
\end{exercise}

\begin{exercise}{457}{曲线约束下的分母简化}
设 $c>0$，$L:y^2=2x$ 从 $(2,2)$ 到 $(2,-2)$。求
$\int_L\frac{(y^2+e^x)\,\dd x+(x^2+y^4)\,\dd y}{y^2-2x+c}$，给出 $c=3$ 的结果；判断 $c=0$ 时能否沿同一路径定义该积分。
\end{exercise}

\begin{exercise}{458}{挖去奇点后的格林公式}
设 $L$ 为 $y^2=x+3,x=2$ 围成区域的逆时针边界。确认原点相对区域的位置，使用挖去小圆的方法计算 $\oint_L\frac{x\,\dd y-y\,\dd x}{x^2+y^2}$；说明仅计算区域内的偏导数差为零为何不足以判定积分为零。
\end{exercise}

\begin{exercise}{459}{移动圆周与奇点穿越}
设 $R>0$，$L_R:(x-1)^2+y^2=R^2$ 取顺时针方向。对 $R<1,R=1,R>1$ 分别讨论
$\oint_{L_R}\frac{y\,\dd x-x\,\dd y}{4x^2+y^2}$ 的存在性并求有定义时的值；可作线性变换把分母化为平方和，须说明方向是否改变。
\end{exercise}

\subsubsection{路径无关、原函数与全微分方程}

\begin{exercise}{460}{参数判定与定义域}
在单位圆盘 $D=\{x^2+y^2<1\}$ 内，求使 $\frac{ax\,\dd x+2y\,\dd y}{x^2+y^2-1}$ 与路径无关的全部 $a$，并求归一化原函数 $u(0,0)=0$。说明单位圆周为何不能包含在该原函数的定义域内。
\end{exercise}

\begin{exercise}{461}{独立函数项的系数匹配}
求使 $(ax^2\cos y-y^3\sin x)\,\dd x+(by^2\cos x+x^3\sin y)\,\dd y$ 在平面内与路径无关的 $a,b$，并找出一个原函数；对所得原函数求偏导数验证，不能只列出系数条件。
\end{exercise}

\begin{exercise}{462}{任意闭曲线积分为零的条件}
设 $F\in C^1(\mathbb R^2)$。证明对任意分段光滑闭曲线 $\Gamma$ 有 $\oint_\Gamma F(x,y)(y\,\dd x+x\,\dd y)=0$ 的充要条件为 $xF_x=yF_y$；分别判断 $F=x+y$ 与 $F=e^{xy}$ 是否满足，并对满足者求一个原函数。
\end{exercise}

\begin{exercise}{463}{分母直线与分支原函数}
求使 $\frac{(ay-x)\,\dd x-y\,\dd y}{(x+y)^2}$ 为全微分的 $a$，并在半平面 $x+y>0$ 内求一个原函数。说明该原函数的表达式如何扩展到 $x+y<0$，以及为何两半平面的积分常数可独立选取。
\end{exercise}

\begin{exercise}{464}{复合原函数与端点公式}
设 $f\in C^1(\mathbb R)$，$L$ 从 $A(a_1,a_2)$ 到 $B(b_1,b_2)$。证明 $\int_Lf(xy)(y\,\dd x+x\,\dd y)$ 与路径无关，并写出仅含一元定积分的端点公式；取 $f(u)=1/(1+u^2)$，$A=(1,0),B=(2,1)$ 计算。
\end{exercise}

\begin{exercise}{465}{由保守场反求未知函数}
设 $f\in C^1(\mathbb R)$，$f(0)=1$，且 $[xy^2-f(x)]\,\dd x+yf(x)\,\dd y$ 与路径无关。求 $f$ 与一个原函数，并计算从 $(0,0)$ 到 $(1,1)$ 的积分；说明初值在唯一确定 $f$ 时的作用。
\end{exercise}

\begin{exercise}{466}{齐次函数与区域条件}
设 $f\in C^1(D)$，$D=\{y>0\}$，且 $f(tx,ty)=t^{-2}f(x,y)$（$t>0$）。证明任意 $D$ 内简单闭曲线 $L$ 均满足 $\oint_Lf(y\,\dd x-x\,\dd y)=0$。再取 $f=1/(x^2+y^2)$，说明把定义域换为去掉原点的整个平面后，同一结论是否仍成立。
\end{exercise}

\begin{exercise}{467}{识别乘积型原函数}
设 $L:x^2/4+y^2/9=1$ 取正向。求 $\oint_L(3x^2\cos x-x^3\sin x)y\,\dd x+(x^3\cos x+2y^2)\,\dd y$，并求从 $(2,0)$ 到 $(0,3)$ 的任意曲线路径积分；必须写出原函数。
\end{exercise}

\begin{exercise}{468}{含非初等项的端点消去}
设 $L:y=\sqrt{2x-x^2}$ 沿上半圆从 $(0,0)$ 到 $(2,0)$，即顺时针方向。证明 $\omega=(e^x+2x\sin y)\,\dd x+(x^2+e^{-y^2})\cos y\,\dd y$ 与路径无关，允许用一元定积分表示原函数，并求 $\int_L\omega$。解释为何计算结果无需求出 $e^{-y^2}\cos y$ 的初等原函数。
\end{exercise}

\begin{exercise}{469}{多项式原函数与端点数据}
设 $L$ 从 $A(1,2)$ 到 $B(-1,1)$。证明 $(x^3-2xy-y^2)\,\dd x-(x+y)^2\,\dd y$ 是全微分并计算其积分；若路径增加一段闭合绕行，说明积分是否变化。
\end{exercise}

\begin{exercise}{470}{径向指数场与路径替换}
设 $L:y=x^3/2$ 从 $(0,0)$ 到 $(2,4)$。计算 $\int_Le^{x^2+y^2}(x\,\dd x+y\,\dd y)$，并求同端点任意分段光滑路径的积分；推广写出 $e^u$ 换为连续函数 $g(u)$ 后的端点公式。
\end{exercise}

\begin{exercise}{471}{线性场的保守性}
证明 $\int_{(-1,-2)}^{(3,4)}(x+2y)\,\dd x+(2x-y)\,\dd y$ 与路径无关并求值。再把第二项中的 $2x$ 改为 $kx$，求仍保持路径无关的 $k$，并说明必要性。
\end{exercise}

\begin{exercise}{472}{混合项配对与原函数}
求 $\int_{(2,0)}^{(4,3)}(xy^2+x-y)\,\dd x+(x^2y-x+y)\,\dd y$。将原函数整理成两个平方项之和，并据此求从 $(0,0)$ 到任意 $(x,y)$ 的积分何时为零。
\end{exercise}

\begin{exercise}{473}{星形线面积与伸缩}
设 $a,b>0$，$L:x=a\cos^3t,y=b\sin^3t$，$0\leq t\leq2\pi$。用曲线积分求围成区域的面积，给出 $a=b$ 时的结果；说明四个尖点是否妨碍按分段光滑边界使用格林公式。
\end{exercise}

\begin{exercise}{474}{双纽线两叶的有向面积}
设 $a>0$，双纽线 $L:(x^2+y^2)^2=a^2(x^2-y^2)$ 的两叶均取正向。用格林公式求总面积；若只把左叶反向，求 $\frac12\int_L(x\,\dd y-y\,\dd x)$，解释它此时是否仍表示总面积。
\end{exercise}

\begin{exercise}{475}{全微分与归一化原函数}
验证 $(2x+\sin y)\,\dd x+x\cos y\,\dd y$ 为全微分，求满足 $u(0,0)=0$ 的原函数，并求其在 $y=0$ 上的最小值；通过偏导数复核原函数。
\end{exercise}

\begin{exercise}{476}{半平面原函数与全局障碍}
求 $\frac{y\,\dd x-x\,\dd y}{x^2+y^2}$ 在右半平面 $x>0$ 内满足 $u(1,0)=0$ 的原函数；判断其能否成为去掉原点的整个平面上的单值连续可微原函数，并用闭圆周积分说明理由。
\end{exercise}

\begin{exercise}{477}{参数匹配与初值归一化}
设 $\dd u=(x+a)e^{x+y}\,\dd x+(xe^{x+y}+bx-2y)\,\dd y$。
\begin{enumerate}[label=(\arabic*)]
\item 求 $a,b$，并说明为何必须在整个平面上匹配偏导数。
\item 已知 $u(0,0)=1$，求 $u(x,y)$。
\item 判断等值线 $u=1$ 在原点附近能否写成 $x=x(y)$，给出依据。
\end{enumerate}
\end{exercise}

\begin{exercise}{478}{全微分方程逐项判定}
分别判断下列方程是否为全微分方程，写出偏导数依据，并求其中全微分方程的隐式通解：
\begin{enumerate}[label=(\arabic*)]
\item $(x^2+y)\,\dd x-x\,\dd y=0$；
\item $(e^y-x)\,\dd x+(y+e^x)\,\dd y=0$；
\item $\sin x\cos y\,\dd x+\cos x\sin y\,\dd y=0$；
\item $xy\,\dd x+(x^2+y^2)\,\dd y=0$。
\end{enumerate}
\end{exercise}

\begin{exercise}{479}{全微分方程与初始等值线}
分别求下列方程的隐式通解，以及经过指定点的等值线方程，并检验所求原函数：
\begin{enumerate}[label=(\arabic*)]
\item $(xy^2-2x+y)\,\dd x-(2y-x-x^2y)\,\dd y=0$，指定点 $(0,1)$；
\item $(x^2+e^{-3y})\,\dd x+(y-3xe^{-3y})\,\dd y=0$，指定点 $(1,0)$。
\end{enumerate}
\end{exercise}

\begin{exercise}{480}{环域上的局部条件与全局原函数}
设 $A=\{1<x^2+y^2<4\}$，
$\omega=\frac{-ay\,\dd x+ax\,\dd y}{x^2+y^2}+b(x\,\dd x+y\,\dd y)$，$a,b\in\mathbb R$。
证明偏导数交叉条件在 $A$ 上始终成立，求任意以原点为圆心的正向圆周（半径在 $(1,2)$ 内）上的积分，据此求 $\omega$ 在 $A$ 内存在单值原函数的充要条件。另求 $1\leq x^2+y^2\leq4$ 的完整正向边界积分，明确内外边界方向。
\end{exercise}

\begin{exercise}{481}{只检验圆周能否判断路径无关}
在去掉原点的平面上，设
\[
 \omega_{a,\lambda}=\frac{-ay\,\dd x+ax\,\dd y}{x^2+y^2}
 +\lambda(y^3\,\dd x+x^3\,\dd y).
\]
求所有以原点为圆心的正向圆周上的积分，确定使这些积分全部为零的参数；再计算正向椭圆 $x^2/4+y^2=1$ 上的积分，求使该形式在整个定义域内与路径无关的全部参数。解释为何圆周检验与最终条件不同。
\end{exercise}

\subsection{曲线积分的应用}

\begin{exercise}{482}{变限积分曲线的弧长}
设 $0<b<\pi/2$，$L_b:y=\int_0^x\tan t\,\dd t$（$0\leq x\leq b$）。求弧长 $s(b)$，给出 $b=\pi/4$ 的结果，并讨论 $b\to\pi/2^-$ 时弧长是否有有限极限。
\end{exercise}

\begin{exercise}{483}{空间螺旋弧的长度与位移}
设 $T>0$，$\Gamma:x=e^t\cos t,y=e^t\sin t,z=e^t$（$0\leq t\leq T$）。求弧长并给出 $T=1$ 的结果；计算两端点距离，证明它严格小于弧长，并说明等号为何不可能成立。
\end{exercise}

\begin{exercise}{484}{零密度点与抛物线细丝}
设 $a>0$，$L_a:y=x^2$（$-a\leq x\leq a$），线密度 $\rho=\sqrt y$。求质量与质心横坐标，给出 $a=1$ 的质量；解释密度在原点为零为何不影响质量严格为正。
\end{exercise}

\begin{exercise}{485}{空间细丝的质量与质心}
设 $L:x=t,y=t^2/2,z=t^3/3$（$0\leq t\leq1$），线密度为 $\sqrt{1+2y+3xz}$。求质量及质心横坐标；指出密度与弧长因子的配合如何把根式积分化为多项式积分。
\end{exercise}

\begin{exercise}{486}{细棒密度与质心位置}
细棒位于 $x$ 轴的 $[0,1]$，密度为 $\rho_c(x)=1+c x(1-x)$，$c\geq0$。求质量与质心，给出 $c=1$ 的结果；用关于中点的对称性解释质心为何不随 $c$ 改变。
\end{exercise}

\begin{exercise}{487}{非均匀螺旋弹簧的重心}
设 $L:x=2\cos t,y=2\sin t,z=3t$（$0\leq t\leq2\pi$），线密度为 $x^2+y^2+z^2$。求质量及全部重心坐标，判断重心高度与几何中点高度 $3\pi$ 的大小关系，并用密度随高度的变化说明结果。
\end{exercise}

\begin{exercise}{488}{半圆铁丝的两个轴惯量}
设 $L:x^2+y^2=1,x\geq0$，线密度 $|y|$。求质量及关于 $x,y$ 两轴的转动惯量，检验两惯量之和与质量的关系；说明半圆区域本身的面积积分为何不能代替本题弧长积分。
\end{exercise}

\begin{exercise}{489}{空间大圆的坐标轴惯量}
设 $a>0$，$\Gamma:x^2+y^2+z^2=a^2,x+y+z=0$，线密度为 $\sqrt{x^2+y^2+z^2}$。求质量与关于三个坐标轴的转动惯量，再求关于过原点且垂直于圆平面的轴的惯量，比较两类轴的结果。
\end{exercise}

\begin{exercise}{490}{圆锥螺线上的变力功}
设 $T>0$，质点从原点沿 $\Gamma_T:x=t\cos t,y=t\sin t,z=t$（$0\leq t\leq T$）在力 $\bm F=(y,-x,z)$ 作用下移动。求功 $W(T)$，给出 $T=\pi/2$ 的结果；求使 $W(T)=0$ 的正 $T$ 并讨论功的符号。
\end{exercise}

\begin{exercise}{491}{中心回复力与椭圆路径}
设 $a,b,k>0$，$\bm F(x,y)=-k(x,y)$。质点沿椭圆 $x^2/a^2+y^2/b^2=1$，从 $(-a,0)$ 顺时针移动至 $(0,b)$。求功并讨论其正负；若沿同椭圆反方向到达相同终点，功是否改变？给出原函数依据。
\end{exercise}

\begin{exercise}{492}{径向斥力与平面距离}
设力方向背离原点，大小为 $k/|z|$，$k>0$。质点沿 $x=at,y=bt,z=ct$（$c\ne0$）从 $t=1$ 移动到 $t=2$。求功，分别核查 $c>0,c<0$ 的表达式；再把终点改为 $t=\lambda>1$，求功关于 $\lambda$ 的关系。
\end{exercise}

\begin{exercise}{493}{细丝惯量的平行轴公式及密度条件}
设分段光滑细丝 $L$ 上的连续线密度 $\rho\geq0$，总质量 $M>0$，质心纵坐标为 $\bar y$。
证明关于直线 $y=c$ 的转动惯量满足
$I(c)=I(\bar y)+M(c-\bar y)^2$，由此证明过质心的水平轴唯一使惯量最小。
若允许有符号密度，取单位圆上 $\rho\equiv-1$，计算 $I(c)$，说明最小化结论为何失效。
\end{exercise}

\begin{exercise}{494}{半圆细丝的质心与旋转力功}
设 $R>0$，上半圆细丝 $L:x^2+y^2=R^2,y\geq0$ 的密度为 $1+\lambda y/R$。
求密度非负的全部 $\lambda$、质量、质心及关于原点的极惯量。再设力场为
$\bm F=(-x-\alpha y,-y+\alpha x)$，$\alpha\in\mathbb R$，求质点沿 $L$ 从 $(R,0)$ 逆时针走到 $(-R,0)$ 的功，讨论功的符号，并解释密度参数是否进入功的表达式。
\end{exercise}

\subsubsection{综合提高题}

\begin{exercise}{495}{边界互换对称性}
设 $L$ 为 $y=x^3$ 与 $y=\sqrt[3]x$ 在第一象限围成区域的完整边界。求 $\oint_L(x-y)^{15}\,\dd s$，并证明将 $15$ 换成任意正奇数结论不变；说明此处不需要指定边界方向的原因。
\end{exercise}

\begin{exercise}{496}{象限分解与两种积分}
设 $L:x^2+y^2=1$，$L_1$ 为第一象限弧，$I=\oint_L(x+y+|x|+|y|)\,\dd s$，$J=\int_{L_1}y\,\dd s$。
\begin{enumerate}[label=(\arabic*)]
\item 用对称性确定 $I,J$ 的关系。
\item 计算两者。
\item 把单位圆换为半径 $R>0$ 的圆，求相应两积分并检验关系是否保留。
\end{enumerate}
\end{exercise}

\begin{exercise}{497}{悬链线上的对数绝对值}
设 $a>0$，$L_a:y=(e^x+e^{-x})/2$（$-a\leq x\leq a$）。求 $\int_{L_a}|\ln(2y-e^x)|\,\dd s$，给出 $a=1$ 的结果；写明对数真数为何为正及去掉绝对值的分段点。
\end{exercise}

\begin{exercise}{498}{空间三角形边界与非负弧长}
设 $L$ 为 $O(0,0,0),A(2,2,0),B(2,2,2)$ 构成的三角形完整边界。分边求 $\oint_Lz\sin(xy)\,\dd s$，指出哪条边贡献为零；将边界方向全部反转后，积分是否改变？须根据弧长微分说明。
\end{exercise}

\begin{exercise}{499}{扇形完整边界的径向积分}
设 $R>0,0<\alpha<\pi$，$L$ 为扇形 $0\leq r\leq R,0\leq\theta\leq\alpha$ 的完整边界。求 $\oint_Le^{\sqrt{x^2+y^2}}\,\dd s$，给出 $R=4,\alpha=\pi/4$ 的结果；分别列出两条半径和圆弧的贡献，讨论漏去半径将产生的误差。
\end{exercise}

\begin{exercise}{500}{等值曲线上的符号判断}
设 $f\in C^1(\mathbb R^2)$，$\Gamma$ 为等值线 $f=1$ 上从第二象限点 $M$ 到第四象限点 $N$ 的非零长度光滑弧。分别判断 $\int_\Gamma f\,\dd x$、$\int_\Gamma f\,\dd y$、$\int_\Gamma f\,\dd s$、$\int_\Gamma(f_x\,\dd x+f_y\,\dd y)$ 的符号并证明；说明这些结论是否需要坐标沿路径单调。
\end{exercise}

\begin{exercise}{501}{曲边闭域的直接积分与补线}
设 $L$ 为 $y=x^2,y^2=x$ 围成区域的顺时针边界。计算 $\oint_L\frac1{\sqrt{1+y}}\,\dd x+x^2\,\dd y$，写出两段方向；另用格林公式写出同值二重积分，核查方向符号。
\end{exercise}

\begin{exercise}{502}{曲线族中的功最小化}
设 $a>0$，$L_a:y=a\sin x$ 从 $(0,0)$ 到 $(\pi,0)$。求 $I(a)=\int_{L_a}(1+y^3)\,\dd x+(2x+y)\,\dd y$，找出使 $I(a)$ 最小的曲线并求最小值；证明该点在全部 $a>0$ 中给出全局最小值。
\end{exercise}

\begin{exercise}{503}{曲线约束与空间闭合积分}
设 $\Gamma$ 为 $x^2+y^2+z^2=4,x^2+y^2=2x,z\geq0$ 的完整上半交线，包含闭合点 $(2,0,0)$，从 $x$ 轴正方向看为逆时针。求
$\oint_\Gamma\frac{y^2\,\dd x+z^2\,\dd y+x^2\,\dd z}{2x+z^2}$。
先用两曲面方程化简分母，参数化时在闭合点处分段或取极限；说明该点是否造成积分发散。
\end{exercise}

\begin{exercise}{504}{第二类积分的估值与极限}
\begin{enumerate}[label=(\arabic*)]
\item 设有向分段光滑曲线 $L$ 长 $K$，$P,Q$ 连续，$M=\max_L\sqrt{P^2+Q^2}$。证明 $|\int_LP\,\dd x+Q\,\dd y|\leq KM$，并说明常向量场沿直线段时何时取等号。
\item 设 $L_R:x^2+y^2=R^2$ 取逆时针方向。用上述不等式估计 $I_R=\oint_{L_R}\frac{y\,\dd x-x\,\dd y}{(x^2+xy+y^2)^2}$，求 $R\to+\infty$ 时的极限，并用缩放说明估计中的 $R$ 次数是否准确。
\end{enumerate}
\end{exercise}

\begin{exercise}{505}{绝对值边界与线性变换}
设 $a>0$，$L_a:|x|+|x+y|=a$ 取逆时针方向。求
$\oint_{L_a}\frac{(x-y)\,\dd x+(x+y)\,\dd y}{|x|+|x+y|}$，给出 $a=1$ 的结果；用 $u=x,v=x+y$ 求包围区域面积，并说明变换是否保持方向。
\end{exercise}

\begin{exercise}{506}{隐函数导数与格林公式}
设 $F\in C^1$，$F(z+x,z-y)=0$ 在单位圆盘邻域内唯一确定 $z=z(x,y)\in C^1$，且沿解有 $F_u+F_v\ne0$。推导 $z_x-z_y$ 的值，计算正向单位圆周 $L$ 上
$\oint_L(z-3y)\,\dd x+(z+x^2)\,\dd y$；说明非退化条件在隐式求导中的作用。
\end{exercise}

\begin{exercise}{507}{正函数环流的下界与等号}
设 $a>0$，$L:(x-a)^2+(y-a)^2=a^2$ 取正向，$f$ 在 $[0,2a]$ 上连续且严格为正。证明
$\oint_Lxf(y)\,\dd y-\frac y{f(x)}\,\dd x\geq2\pi a^2$，并求等号成立的充要条件。证明须说明为何只假设 $f$ 连续即可使用所需的偏导数。
\end{exercise}

\begin{exercise}{508}{余弦弧补线与奇偶性}
设 $L:y=\cos x$ 从 $(-\pi/2,0)$ 到 $(\pi/2,0)$。求
$\int_L(e^x-y^3+xy^2)\,\dd x-(x^2y+2x+\sin y)\,\dd y$。
将结果分为端点项和补线后的面积积分项，并说明奇函数项在哪一步消去。
\end{exercise}

\begin{exercise}{509}{最优区域与带奇点环流}
在平面有界单连通闭区域中，令 $D_1$ 为使 $I(D)=\iint_D(4-x^2-y^2)\,\dd x\dd y$ 最大的闭圆盘。
\begin{enumerate}[label=(\arabic*)]
\item 求 $D_1$ 与 $I(D_1)$，证明任何额外包含负值区域或遗漏正值区域都不能改善积分。
\item 计算正向边界上的
$\oint_{\partial D_1}\frac{(xe^{x^2+4y^2}+y)\,\dd x+(4ye^{x^2+4y^2}-x)\,\dd y}{x^2+4y^2}$；区分可写为全微分的部分与角形式部分。
\end{enumerate}
\end{exercise}

\begin{exercise}{510}{路径无关诱导的微分方程}
设 $f\in C^1(\mathbb R)$，$f(0)=0$，且 $[f(x)-e^x]\sin y\,\dd x-f(x)\cos y\,\dd y$ 与路径无关。求 $f$，再求一个原函数并求从 $(0,0)$ 到 $(1,\pi/2)$ 的积分。
\end{exercise}

\begin{exercise}{511}{同圆不同弧的端点一致性}
设 $L$ 为圆周 $x^2+y^2=2y$ 上从 $O(0,0)$ 到 $A(1,1)$ 的任一段有向弧。求 $\int_L(1-2xy-y^2)\,\dd x-(x+y)^2\,\dd y$，证明选择较长或较短的圆弧结果相同，并写出原函数。
\end{exercise}

\begin{exercise}{512}{右半平面中的反求函数}
设 $f$ 在 $x>0$ 可导，$f(1)=0$，$yf(x)\,\dd x+[2xf(x)-x^2]\,\dd y$ 在右半平面内与路径无关。
\begin{enumerate}[label=(\arabic*)]
\item 求 $f(x)$，写出由偏导数条件得到的一阶微分方程。
\item 求从 $(1,0)$ 到 $(2,3)$ 的积分，并验证所求原函数。
\end{enumerate}
\end{exercise}

\begin{exercise}{513}{端点恒等式反求场分量}
设 $Q\in C^1(\mathbb R^2)$，$2xy\,\dd x+Q(x,y)\,\dd y$ 与路径无关，且从原点到 $(t,1)$ 与到 $(1,t)$ 的积分对所有实数 $t$ 相等。求 $Q$ 及满足 $u(0,0)=0$ 的原函数，说明“对所有 $t$”如何确定任意函数项。
\end{exercise}

\begin{exercise}{514}{偏导信息、端点积分与最小值}
设 $f_x=(2x+1)e^{2x-y}$，$f(0,y)=y+1$，$L_t$ 从 $(0,0)$ 到 $(1,t)$。
\begin{enumerate}[label=(\arabic*)]
\item 求 $f(x,y)$，核查边界条件。
\item 求 $I(t)=\int_{L_t}f_x\,\dd x+f_y\,\dd y$。
\item 在全部 $t\in\mathbb R$ 中求 $I(t)$ 的最小值，并证明最小点唯一。
\end{enumerate}
\end{exercise}

\begin{exercise}{515}{三角形有向面积公式}
设三角形顶点按逆时针顺序为 $M_i(x_i,y_i)$（$i=1,2,3$）。用曲线积分证明
$S=\frac12[(x_1y_2-x_2y_1)+(x_2y_3-x_3y_2)+(x_3y_1-x_1y_3)]$。
再说明顶点顺序反转及三个点共线时该表达式的含义，给出不预设方向时的面积公式。
\end{exercise}

\begin{exercise}{516}{全微分参数与等值线方程}
设 $\dd u=(xy^2+a\sin x\cos x\sin3y)\,\dd x-(bx^2y+\cos2x\cos3y)\,\dd y$。
\begin{enumerate}[label=(\arabic*)]
\item 求 $a,b$。
\item 求全部原函数 $u$，并给出 $u(0,0)=0$ 的归一化结果。
\item 求将该微分形式置零所得方程的隐式通解，并求经过 $(0,0)$ 的等值线。
\end{enumerate}
\end{exercise}

\begin{exercise}{517}{不同曲线的弧长比较}
设 $s_1$ 为 $y=\sin x$（$0\leq x\leq2\pi$）的弧长，$s_2$ 为椭圆 $x^2+2y^2=2$ 的周长。写出两者的定积分表达式，通过变量代换确定其关系；不要求计算非初等积分的具体值，并说明比较为何不能仅靠图形外观。
\end{exercise}

\begin{exercise}{518}{微分方程曲线的均匀重心}
设 $y(x)$ 满足 $2xy'-4y=2\ln x-1$、$y(1)=1/4$。求 $y(x)$，化简弧长因子，求均匀曲线弧 $1\leq x\leq e$ 的质量（单位线密度）及重心的两个坐标；用 $1<\bar x<e$ 检验结果。
\end{exercise}

\begin{exercise}{519}{旋转径向力与角度功}
质点在 $(x,y)\ne(0,0)$ 受到力，大小为到原点距离的倒数的两倍，方向为 $\overrightarrow{OM}$ 顺时针旋转 $90^\circ$ 的方向。写出力的坐标分量，求从 $A(0,\pi/2)$ 沿 $y=(\pi/2)\cos x$ 到 $B(\pi/2,0)$ 的功；再判断换成第一象限内同端点的任意路径，功是否改变，并说明定义域条件。
\end{exercise}
